\begin{frontmatter}

\title{Cheap Signals, Costly Proof: Award-Layer Evidence Triage in Cartel Enforcement\texorpdfstring{\footnote{We thank seminar participants at INSPER for comments; all remaining errors are our own.}}{}}

\author[insper]{Darcio Genicolo-Martins\corref{cor1}}
\ead{darciogm1@insper.edu.br}
\author[insper]{Paulo Furquim de Azevedo}
\affiliation[insper]{organization={INSPER},
  city={S\~ao Paulo}, country={Brazil}}
\cortext[cor1]{Corresponding author.}

\begin{abstract}
Cartel enforcement begins before cartel proof exists. Agencies
often observe contract-award records long before they recover
the bid-level evidence needed to prove coordination. This paper
studies whether that cheap record layer can allocate costly
forensic attention. Using S\~ao Paulo's BEC procurement platform
($\valYearStart$--$\valYearEnd$) and CADE adjudications, we
construct a \emph{frequent-loser} screen from award records
alone: among zero-win firms, persistent participation ranks
loser-side cartel-adjacency risk. The screen targets forensic
priority: it orders where proof-producing bid analysis should
begin. We validate the ranking against adjudication-anchored
cobidder labels, discipline it with timing and exposure audits,
and compare it with bid-distribution forensics on the same
target. The award-layer signal adds non-redundant information to
the bid layer. Used sequentially, the award-layer ranking
reduces the bid-microdata pool for the forensic stage by 83\%
while recovering 131 of 193 adjudicated cobidders. The results
support an enforcement-design claim: cheap administrative
records can discipline where costly proof-producing
investigation should begin while leaving liability to the richer
evidentiary record.
\end{abstract}

\end{frontmatter}

\noindent\textbf{Keywords:} cartel screening, public procurement, award-layer data, enforcement triage, incomplete observability, bid rigging

\noindent\textbf{JEL No.:} D44, D73, H57, K21, L41

\bigskip
